\newpage
\section{Motivação}
\label{sc:motivacao}

A intensificação dos eventos climáticos extremos e o aumento expressivo da ocorrência de queimadas e incêndios florestais no território brasileiro configuram uma das ameaças mais severas à segurança socioambiental e à estabilidade das infraestruturas críticas nacionais. Dentre essas infraestruturas, destaca-se o Sistema Interligado Nacional (SIN), cuja malha de transmissão de energia elétrica abrange mais de 2.346 linhas de alta e extra-alta tensão ($\ge 230\text{ kV}$), estendendo-se por mais de 182 mil quilômetros e interligando os subsistemas Norte, Nordeste, Sudeste/Centro-Oeste e Sul.

Historicamente, as queimadas sob faixas de servidão constituem uma das principais causas de desligamentos forçados de linhas de transmissão no Brasil. O calor extremo emitido pelas chamas e o acúmulo de partículas condutoras de fuligem e fuligem ionizada presentes na pluma de fumaça reduzem drasticamente a rigidez dielétrica do ar ao redor dos condutores, provocando descargas elétricas disruptivas fase-terra conhecidas tecnicamente como \textit{flashover}. Esse fenômeno resulta na atuação imediata dos sistemas de proteção e no desligamento intempestivo da linha, gerando desequilíbrio transitório na rede e, em situações críticas, cortes maciços de suprimento de energia para as distribuidoras e centros urbanos.

Dados abertos apurados pelo Operador Nacional do Sistema Elétrico (ONS) \cite{ons2026dadosabertos} revelam que, entre janeiro de 2007 e novembro de 2026, foram contabilizados mais de 8.395 eventos de corte severo de carga (perturbações $\ge 100\text{ MW}$ por $\ge 10\text{ minutos}$ conforme os Procedimentos de Rede da norma RO-AO.BR.03), totalizando um montante acumulado superior a 717.448 MW de carga interrompida. Essa fragilidade é amplificada pela marcante sazonalidade climática: os meses do trimestre seco (agosto, setembro e outubro) concentram mais de 35\% de todas as ocorrências de desligamento forçado do país, coincidindo com o período de estiagem crítica no bioma Cerrado e na Amazônia Legal, onde a umidade relativa do ar atinge valores inferiores a 15\% e as rajadas de vento superam 45 km/h.

Apesar da magnitude do problema e da disponibilidade de dados geoespaciais abertos provenientes de agências espaciais e institutos meteorológicos (como o Instituto Nacional de Pesquisas Espaciais - INPE \cite{inpe2025bdqueimadas}, o Instituto Nacional de Meteorologia - INMET e a NASA), as ferramentas convencionais adotadas por operadores de centros de controle e agentes da Defesa Civil apresentam limitações estruturais:

\begin{enumerate}
    \item \textbf{Passividade dos Painéis de Monitoramento (Dashboards)}: Os centros de operação dependem quase exclusivamente de telas analíticas passivas. Em situações de conflagração generalizada (como a crise de queimadas registrada em agosto e setembro de 2024 no estado de São Paulo e no Pantanal), o operador é sobrecarregado por centenas de pontos luminosos no mapa, sem priorização dinâmica nem despacho automático de alertas.
    \item \textbf{Isotropia Irreal dos Modelos de Buffer}: A maioria dos sistemas existentes projeta círculos concêntricos equidistantes (buffers isotrópicos) ao redor dos focos de calor detectados. Sob a ação de ventos intensos, o fogo não se propaga em círculo, mas em elipses altamente orientadas a sotavento, superestimando o risco a barlavento e negligenciando a aproximação veloz da frente de chama contra condutores elétricos vulneráveis.
    \item \textbf{Latência na Geração de Relatórios Situacionais}: A confecção de relatórios de situação (SitRep) para órgãos reguladores e brigadistas costuma ser um processo manual, demandando horas de triagem entre diferentes fontes governamentais e atrasando a tomada de decisão em campo.
\end{enumerate}

Nesse contexto, torna-se imperativa a pesquisa e o desenvolvimento de uma arquitetura DataOps geoespacial orientada a eventos que transcenda a mera visualização passiva. A integração sinérgica entre sensoriamento orbital multissensor (satélites LEO e GEO), modelagem física direcional ancorada em vetores meteorológicos, indexação espacial em grades hexagonais discretas (Uber H3) e agentes de Inteligência Artificial para geração instantânea de boletins executivos fundamenta este Trabalho de Graduação, preenchendo uma lacuna crítica entre a ciência de dados e a resposta operacional a desastres no Brasil.
